\section{Atomic Interfaces}
\label{sec:atomic}

Atomicity gives the simplest sufficient condition for reportability.  If each
active contract atom has exactly one active lever, then every subset of the
active lever interface is block-aligned.  Pairwise block disjointness yields
the complement identity, and residual faithfulness transfers feasibility and
minimality between the implementation and contract domains.

\begin{theorem}[M3-A: atomic grouped correctness]
\label{thm:m3a}
Assume active blocks are pairwise disjoint.  If \(\RepairAtomicity\) and
\(\ResidualFaithfulness\) hold, then for every contract deletion \(G\),
\[
  \GroupedRepair(G)
  \quad\Longleftrightarrow\quad
  \ContractRepair(G).
\]
\end{theorem}

The Lean declaration is \code{m3a_grouped_correctness}.  The proof factors
through block-alignment closure, complement correspondence, feasibility
transfer, minimality transfer, and the identity
\(\groupTouchAny(\betaSet(G))=G\).

\begin{corollary}[Atom-indexed raw transport]
\label{cor:m3a-raw}
For two atomic, residually faithful realizations of the same contract, and for
every \(G\subseteq I\), the block deletion induced by \(G\) is a raw repair in
the first realization if and only if the block deletion induced by \(G\) is a
raw repair in the second.
\end{corollary}

This is the Lean theorem \code{m3a_raw_transport}.  It is atom-indexed rather
than an arbitrary lever-bijection theorem.

\paragraph{Limit.}
Atomicity is sufficient, not necessary.  M3-A does not apply to Candidate B
because
\[
  \beta(\mathtt{minlib})
  =
  \{\mathtt{decisive\_voter0},\mathtt{decisive\_voter1}\}
\]
is non-atomic.  Raw transport is an M3-A result only and must not be attributed
to M3-B.

\paragraph{Dependency record.}
Appendix~\ref{app:lean} includes the exact theorem signatures.  The longer
atomicity-use dependency table will be placed with the proof details in the
technical appendix during full drafting.
